% GENERATED FILE. Edit canonical lessons, then run npm run generate:books.
% canonical-source-hash: fnv1a64:223863b244180d4d
% canonical-lessons: PA-C245-vihra, PA-C245-kotha, PA-C245-chubara, PA-C245-thammh, PA-C245-itt, PA-R245-review-239-245

\chapter{Our Courtyard House}
\label{ch:pa-a2-our-courtyard-house}
\hlchaptermodality{245}

\begin{chapteropening}
\textbf{By the end of this chapter:} \emph{I can name the parts of a courtyard house.}

Complete the last lesson of chapter 245: I can name the parts of a courtyard house.
\end{chapteropening}

\section[\texorpdfstring{\pa{ਵਿਹੜਾ} (vihṛā)}{ਵਿਹੜਾ (vihṛā)}]{\pa{ਵਿਹੜਾ} (vihṛā) — a courtyard}

\label{lesson:PA-C245-vihra}



\begin{quote}
Before the new one: say the Punjabi for to serve food, then the Punjabi for to mash.
\end{quote}

\subsection*{You'll want to know: \pa{ਵਿਹੜਾ}}
\textbf{\pa{ਵਿਹੜਾ}} — \emph{vihṛā} — \textquotedblleft{}a courtyard\textquotedblright{}.

\begin{quote}
\pa{ਸਾਡੇ} \pa{ਘਰ} \pa{ਦਾ} \pa{ਵਿਹੜਾ} \pa{ਵੱਡਾ} \pa{ਹੈ।}
\end{quote}

\begin{quote}
\emph{sāḍe ghar dā vihṛā vaḍḍā hai.}
\end{quote}

\begin{quote}
Our house has a big courtyard.
\end{quote}

A masculine noun.

\begin{cousinweb}
\textbf{\pa{ਵਿਹੜਾ}} (\emph{vihṛā}) is inherited from Sanskrit \emph{vēṣṭa}, \textquotedblleft{}band, enclosure\textquotedblright{}; CDIAL lists Punjabi \emph{veṛhā}, \textquotedblleft{}enclosure, courtyard\textquotedblright{} (CDIAL 12130).
\end{cousinweb}

\begin{grammarlens}[title={What you've built}]
One more word for this chapter.
\end{grammarlens}

\subsection*{Your turn}
\begin{itemize}
\raggedright
  \item \emph{Say it:} \emph{vihṛā}
  \item \emph{Say it:} \emph{vihṛā}, once more
  \item \emph{Say it:} \emph{sāḍe ghar dā vihṛā vaḍḍā hai.} — Our house has a big courtyard.
  \item \emph{Recall:} say the Punjabi for to serve food, then the Punjabi for to mash, then say \emph{vihṛā} again
\end{itemize}

\begin{tcolorbox}[breakable,colback=teal!4,colframe=teal!35!black,title={Before you move on}]
What does \pa{ਵਿਹੜਾ} mean? (\textquotedblleft{}A courtyard\textquotedblright{}.) Say it once more.
\end{tcolorbox}
\section[\texorpdfstring{\pa{ਕੋਠਾ} (koṭhā)}{ਕੋਠਾ (koṭhā)}]{\pa{ਕੋਠਾ} (koṭhā) — a flat-roofed house; a room}

\label{lesson:PA-C245-kotha}



\begin{quote}
Before the new one: say the Punjabi for to mash, then the Punjabi for a courtyard.
\end{quote}

\subsection*{You'll want to know: \pa{ਕੋਠਾ}}
\textbf{\pa{ਕੋਠਾ}} — \emph{koṭhā} — \textquotedblleft{}a flat-roofed house; a room\textquotedblright{}.

\begin{quote}
\pa{ਪਿੰਡ} \pa{ਵਿੱਚ} \pa{ਦਾਦਾ} \pa{ਜੀ} \pa{ਦਾ} \pa{ਕੋਠਾ} \pa{ਪੁਰਾਣਾ} \pa{ਹੈ।}
\end{quote}

\begin{quote}
\emph{piṇḍ vich dādā jī dā koṭhā purāṇā hai.}
\end{quote}

\begin{quote}
Grandfather's house in the village is old.
\end{quote}

A masculine noun.

\begin{cousinweb}
\textbf{\pa{ਕੋਠਾ}} (\emph{koṭhā}) is inherited from Sanskrit \emph{kōṣṭha}, \textquotedblleft{}granary, storeroom, inner apartment\textquotedblright{}; CDIAL lists Punjabi \emph{koṭhā}, \textquotedblleft{}house with mud roof and walls\textquotedblright{} (CDIAL 3546).
\end{cousinweb}

\begin{grammarlens}[title={What you've built}]
One more word for this chapter.
\end{grammarlens}

\subsection*{Your turn}
\begin{itemize}
\raggedright
  \item \emph{Say it:} \emph{koṭhā}
  \item \emph{Say it:} \emph{koṭhā}, once more
  \item \emph{Say it:} \emph{piṇḍ vich dādā jī dā koṭhā purāṇā hai.} — Grandfather's house in the village is old.
  \item \emph{Recall:} say the Punjabi for to mash, then the Punjabi for a courtyard, then say \emph{koṭhā} again
  \item \emph{Write it:} copy \textbf{\pa{ਕੋਠਾ}} once, checking the whole word
\end{itemize}

\begin{tcolorbox}[breakable,colback=teal!4,colframe=teal!35!black,title={Before you move on}]
What does \pa{ਕੋਠਾ} mean? (\textquotedblleft{}A flat-roofed house; a room\textquotedblright{}.) Say it once more.
\end{tcolorbox}
\section[\texorpdfstring{\pa{ਚੁਬਾਰਾ} (chubārā)}{ਚੁਬਾਰਾ (chubārā)}]{\pa{ਚੁਬਾਰਾ} (chubārā) — an upper room}

\label{lesson:PA-C245-chubara}



\begin{quote}
Before the new one: say the Punjabi for a courtyard, then the Punjabi for a flat-roofed house.
\end{quote}

\subsection*{You'll want to know: \pa{ਚੁਬਾਰਾ}}
\textbf{\pa{ਚੁਬਾਰਾ}} — \emph{chubārā} — \textquotedblleft{}an upper room\textquotedblright{}.

\begin{quote}
\pa{ਚੁਬਾਰੇ} \pa{ਵਿੱਚ} \pa{ਚਾਰ} \pa{ਖਿੜਕੀਆਂ} \pa{ਹਨ।}
\end{quote}

\begin{quote}
\emph{chubāre vich chār khiṛkīāṁ han.}
\end{quote}

\begin{quote}
The upper room has four windows.
\end{quote}

A masculine noun.

\begin{cousinweb}
\textbf{\pa{ਚੁਬਾਰਾ}} (\emph{chubārā}) is inherited from Sanskrit \emph{caturdvāra}, \textquotedblleft{}having four doors\textquotedblright{}; CDIAL lists Punjabi \emph{cubārā}, \textquotedblleft{}upper room with four windows\textquotedblright{} (CDIAL 4611).
\end{cousinweb}

\begin{grammarlens}[title={What you've built}]
One more word for this chapter.
\end{grammarlens}

\subsection*{Your turn}
\begin{itemize}
\raggedright
  \item \emph{Say it:} \emph{chubārā}
  \item \emph{Say it:} \emph{chubārā}, once more
  \item \emph{Say it:} \emph{chubāre vich chār khiṛkīāṁ han.} — The upper room has four windows.
  \item \emph{Recall:} say the Punjabi for a courtyard, then the Punjabi for a flat-roofed house, then say \emph{chubārā} again
  \item \emph{Write it:} copy \textbf{\pa{ਚੁਬਾਰਾ}} once, checking the whole word
\end{itemize}

\begin{tcolorbox}[breakable,colback=teal!4,colframe=teal!35!black,title={Before you move on}]
What does \pa{ਚੁਬਾਰਾ} mean? (\textquotedblleft{}An upper room\textquotedblright{}.) Say it once more.
\end{tcolorbox}
\section[\texorpdfstring{\pa{ਥੰਮ੍ਹ} (thammh)}{ਥੰਮ੍ਹ (thammh)}]{\pa{ਥੰਮ੍ਹ} (thammh) — a pillar}

\label{lesson:PA-C245-thammh}



\begin{quote}
Before the new one: say the Punjabi for a flat-roofed house, then the Punjabi for an upper room.
\end{quote}

\subsection*{You'll want to know: \pa{ਥੰਮ੍ਹ}}
\textbf{\pa{ਥੰਮ੍ਹ}} — \emph{thammh} — \textquotedblleft{}a pillar\textquotedblright{}.

\begin{quote}
\pa{ਵਿਹੜੇ} \pa{ਵਿੱਚ} \pa{ਦੋ} \pa{ਥੰਮ੍ਹ} \pa{ਹਨ।}
\end{quote}

\begin{quote}
\emph{vihṛe vich do thammh han.}
\end{quote}

\begin{quote}
There are two pillars in the courtyard.
\end{quote}

A masculine noun.

\begin{cousinweb}
\textbf{\pa{ਥੰਮ੍ਹ}} (\emph{thammh}) is inherited from Sanskrit \emph{stambha}, \textquotedblleft{}pillar, post\textquotedblright{}; CDIAL lists Punjabi \emph{thammh}, \textquotedblleft{}pillar, post\textquotedblright{} (CDIAL 13682).
\end{cousinweb}

\begin{grammarlens}[title={What you've built}]
One more word for this chapter.
\end{grammarlens}

\subsection*{Your turn}
\begin{itemize}
\raggedright
  \item \emph{Say it:} \emph{thammh}
  \item \emph{Say it:} \emph{thammh}, once more
  \item \emph{Say it:} \emph{vihṛe vich do thammh han.} — There are two pillars in the courtyard.
  \item \emph{Recall:} say the Punjabi for a flat-roofed house, then the Punjabi for an upper room, then say \emph{thammh} again
  \item \emph{Write it:} copy \textbf{\pa{ਥੰਮ੍ਹ}} once, checking the whole word
\end{itemize}

\begin{tcolorbox}[breakable,colback=teal!4,colframe=teal!35!black,title={Before you move on}]
What does \pa{ਥੰਮ੍ਹ} mean? (\textquotedblleft{}A pillar\textquotedblright{}.) Say it once more.
\end{tcolorbox}
\section[\texorpdfstring{\pa{ਇੱਟ} (iṭṭ)}{ਇੱਟ (iṭṭ)}]{\pa{ਇੱਟ} (iṭṭ) — a brick}

\label{lesson:PA-C245-itt}



\begin{quote}
Before the new one: say the Punjabi for an upper room, then the Punjabi for a pillar.
\end{quote}

\subsection*{You'll want to know: \pa{ਇੱਟ}}
\textbf{\pa{ਇੱਟ}} — \emph{iṭṭ} — \textquotedblleft{}a brick\textquotedblright{}.

\begin{quote}
\pa{ਕੰਧ} \pa{ਲਾਲ} \pa{ਇੱਟਾਂ} \pa{ਦੀ} \pa{ਹੈ।}
\end{quote}

\begin{quote}
\emph{kandh lāl iṭṭāṁ dī hai.}
\end{quote}

\begin{quote}
The wall is made of red bricks.
\end{quote}

A feminine noun.

\begin{cousinweb}
\textbf{\pa{ਇੱਟ}} (\emph{iṭṭ}) is inherited from Sanskrit \emph{iṣṭakā}, \textquotedblleft{}brick\textquotedblright{}; CDIAL lists Punjabi \emph{iṭṭ}, \textquotedblleft{}brick\textquotedblright{} (CDIAL 1600).
\end{cousinweb}

\begin{grammarlens}[title={What you've built}]
One more word for this chapter.
\end{grammarlens}

\subsection*{Your turn}
\begin{itemize}
\raggedright
  \item \emph{Say it:} \emph{iṭṭ}
  \item \emph{Say it:} \emph{iṭṭ}, once more
  \item \emph{Say it:} \emph{kandh lāl iṭṭāṁ dī hai.} — The wall is made of red bricks.
  \item \emph{Recall:} say the Punjabi for an upper room, then the Punjabi for a pillar, then say \emph{iṭṭ} again
  \item \emph{Write it:} copy \textbf{\pa{ਇੱਟ}} once, checking the whole word
\end{itemize}

\begin{tcolorbox}[breakable,colback=teal!4,colframe=teal!35!black,title={Before you move on}]
What does \pa{ਇੱਟ} mean? (\textquotedblleft{}A brick\textquotedblright{}.) Say it once more.
\end{tcolorbox}
\section[(dialogue)]{Looking back: from Bread on the Griddle to Our Courtyard House}

\label{lesson:PA-R245-review-239-245}



\begin{quote}
Say the Punjabi for a griddle and to bake on a griddle.
\end{quote}

\subsection*{Your turn}
Say these aloud:

\begin{itemize}
\raggedright
  \item a griddle, to bake on a griddle, tongs, to cover, ghee
  \item a grinding stone, to grind, a wooden mortar, a pestle, to pound
  \item to peel, to chop finely, to split, a chilli, gram flour
  \item to fry, to roast, a ladle, to sprinkle, to get burnt
  \item to pour in, to dissolve, to mix, to sift, to taste
  \item a straining cloth, to strain, to pour off, to serve food, to mash
  \item a courtyard, a flat-roofed house, an upper room, a pillar, a brick
\end{itemize}

\begin{tcolorbox}[breakable,colback=teal!4,colframe=teal!35!black,title={Before you move on}]
A griddle? (\textbf{\pa{ਤਵਾ}}, \emph{tavā}.) To bake on a griddle? (\textbf{\pa{ਸੇਕਣਾ}}, \emph{sekṇā}.) Tongs? (\textbf{\pa{ਚਿਮਟਾ}}, \emph{chimṭā}.) To cover? (\textbf{\pa{ਢੱਕਣਾ}}, \emph{ḍhakkṇā}.) Ghee? (\textbf{\pa{ਘਿਉ}}, \emph{ghiu}.) A grinding stone? (\textbf{\pa{ਸਿਲ}}, \emph{sil}.) To grind? (\textbf{\pa{ਪੀਹਣਾ}}, \emph{pīhṇā}.) A wooden mortar? (\textbf{\pa{ਉੱਖਲੀ}}, \emph{ukkhlī}.) A pestle? (\textbf{\pa{ਮੋਹਲਾ}}, \emph{mohlā}.) To pound? (\textbf{\pa{ਕੁੱਟਣਾ}}, \emph{kuṭṭṇā}.) To peel? (\textbf{\pa{ਛਿੱਲਣਾ}}, \emph{chhillṇā}.) To chop finely? (\textbf{\pa{ਕਤਰਨਾ}}, \emph{katarnā}.) To split? (\textbf{\pa{ਚੀਰਨਾ}}, \emph{chīrnā}.) A chilli? (\textbf{\pa{ਮਿਰਚ}}, \emph{mirach}.) Gram flour? (\textbf{\pa{ਵੇਸਣ}}, \emph{vesaṇ}.) To fry? (\textbf{\pa{ਤਲਣਾ}}, \emph{talṇā}.) To roast? (\textbf{\pa{ਭੁੰਨਣਾ}}, \emph{bhunnṇā}.) A ladle? (\textbf{\pa{ਕੜਛੀ}}, \emph{kaṛchhī}.) To sprinkle? (\textbf{\pa{ਛਿੜਕਣਾ}}, \emph{chhiṛakṇā}.) To get burnt? (\textbf{\pa{ਸੜਨਾ}}, \emph{saṛnā}.) To pour in? (\textbf{\pa{ਪਾਉਣਾ}}, \emph{pāuṇā}.) To dissolve? (\textbf{\pa{ਘੋਲਣਾ}}, \emph{gholṇā}.) To mix? (\textbf{\pa{ਰਲਾਉਣਾ}}, \emph{ralāuṇā}.) To sift? (\textbf{\pa{ਛਾਣਨਾ}}, \emph{chhāṇnā}.) To taste? (\textbf{\pa{ਚੱਖਣਾ}}, \emph{chakkhṇā}.) A straining cloth? (\textbf{\pa{ਪੋਣਾ}}, \emph{poṇā}.) To strain? (\textbf{\pa{ਪੁਣਨਾ}}, \emph{puṇnā}.) To pour off? (\textbf{\pa{ਨਿਤਾਰਨਾ}}, \emph{nitārnā}.) To serve food? (\textbf{\pa{ਪਰੋਸਣਾ}}, \emph{parosṇā}.) To mash? (\textbf{\pa{ਘੋਟਣਾ}}, \emph{ghoṭṇā}.) A courtyard? (\textbf{\pa{ਵਿਹੜਾ}}, \emph{vihṛā}.) A flat-roofed house? (\textbf{\pa{ਕੋਠਾ}}, \emph{koṭhā}.) An upper room? (\textbf{\pa{ਚੁਬਾਰਾ}}, \emph{chubārā}.) A pillar? (\textbf{\pa{ਥੰਮ੍ਹ}}, \emph{thammh}.) A brick? (\textbf{\pa{ਇੱਟ}}, \emph{iṭṭ}.)
\end{tcolorbox}
